\section{Support-coordinate intervention table}
\label{app:support_coordinate_interventions}

\Cref{fig:support_coordinate_interventions} plots the full curves for one Local-linear gates LISTA checkpoint at seed \(0\). \Cref{tab:support_coordinate_interventions_h21} reports accumulated MSE at \(H=21\), defined here as the sum over rollout steps of the squared error averaged across state coordinates. This differs from the state-summed, through-horizon mean used by the headline forecasting table, so their numerical scales are not comparable. Candidate starts lie in the tie-inclusive per-native-label high center-margin subset and must retain their initial native label throughout the true \(21\)-step future. The evaluator balances \(100\) starts across the three labels (\(34/33/33\)). Coordinate dropping and random support use exactly these same starts; random support uses \(20\) inactive-coordinate reassignments per start. Random-support summaries pool the resulting \(2{,}000\) dependent point--shuffle outcomes at each horizon. The margin is not a true basin-boundary distance. This one-system, one-model, one-seed experiment is descriptive and has no confirmatory significance test.

\begin{table}[tbp]
\centering
\caption{Accumulated MSE (sum of per-coordinate step MSE) at \(H=21\) for descriptive support-coordinate interventions on the Local-linear gates LISTA checkpoint at seed \(0\). Drop rows summarize \(100\) starts; the random row pools \(2{,}000\) dependent point--shuffle outcomes. Lower is better; no confirmatory inference is attached.}
\label{tab:support_coordinate_interventions_h21}
\resizebox{0.7\textwidth}{!}{\begin{tabular}{@{}lcc@{}}
\toprule
Rollout & Mean \(\pm\) SD & Median [IQR] \\
\midrule
Standard & $0.0158{\pm}0.0127$ & $0.0140\,[0.0055,0.0224]$ \\
Drop top-1 & $0.508{\pm}0.647$ & $0.218\,[0.148,0.677]$ \\
Drop top-2 & $1.37{\pm}0.938$ & $1.06\,[0.802,1.63]$ \\
Drop top-3 & $2.08{\pm}1.10$ & $1.80\,[1.34,2.10]$ \\
Drop top-5 & $3.26{\pm}2.44$ & $1.95\,[1.61,4.17]$ \\
Drop top-10 & $8.51{\pm}6.83$ & $5.16\,[4.58,8.11]$ \\
Random support & $1.69{\times}10^{3}{\pm}2.19{\times}10^{3}$ & $874\,[377,2.06{\times}10^{3}]$ \\
\bottomrule
\end{tabular}
}
\end{table}
